%\chapter{Cài đặt và thử nghiệm}
\chapter{Implementation and Experiment}
To evaluate the practical capabilities and performance of the Segment Anything Model (SAM), this study utilizes the official implementation provided by Meta AI Research. The core architecture remains consistent with the original proposal; however, the codebase was streamlined by removing extraneous files and modules not required for this specific demonstration. The complete source code and configuration files used in these experiments are available at the following repository: \url{https://github.com/facebookresearch/segment-anything}.

\section{Experimental Environment}
The experiments were conducted in a \textbf{Python 3.11} environment. To ensure the smooth execution of the Segment Anything Model and the evaluation of the results, the following core libraries and frameworks were utilized:

\begin{itemize}
	\item \textbf{Deep Learning Frameworks:} \texttt{torch} and \texttt{torchvision} for model loading and inference.
	\item \textbf{Optimization and Deployment:} \texttt{onnx}, \texttt{onnxruntime}, and \texttt{onnxscript} for model exportation and runtime acceleration.
	\item \textbf{Image Processing and Data Handling:} \texttt{opencv-python} for real-time computer vision tasks and \texttt{pycocotools} for performance metric calculations.
	\item \textbf{Visualization and Development:} \texttt{matplotlib} for generating result plots and \texttt{jupyter} for interactive experimentation.
\end{itemize}

\section{Model Checkpoints and Demonstration Modes}
Due to storage constraints, model weights are not included directly in the repository. Users are encouraged to download the pre-trained checkpoints as specified in the \texttt{Source\_code/README.md} file. The implementation supports three distinct Vision Transformer (ViT) backbones: \textbf{SAM-H} (\texttt{vit\_h}), \textbf{SAM-L} (\texttt{vit\_l}), and \textbf{SAM-B} (\texttt{vit\_b}).

The demonstration is structured into three primary interfaces:

\begin{enumerate}
	\item \textbf{Interactive Notebooks:} Located in \texttt{Source\_code/notebooks/}, these provide a guided environment for manual prompting and automatic mask generation. Users can follow the inline instructions within \texttt{predictor\_example.ipynb} and\\ \texttt{automatic\_mask\_generator\_example.ipynb}.
	
	\item \textbf{Command Line Interface (CLI):} For batch processing, a Python script is provided. After navigating to the \texttt{Source\_code} directory, the following command structure is used:
	\begin{lstlisting}
		python scripts/amg.py --checkpoint <path/to/checkpoint> --model-type <model_type> --input <image_or_folder> --output <path/to/output>
	\end{lstlisting}
	
	\item \textbf{Web-Based Application:} A specialized demo requiring \textbf{Node.js} and \textbf{Yarn}. This mode involves a two-step process: exporting image embeddings and the ONNX model before building the static web application. Detailed setup steps are found in \texttt{Source\_code/demo/README.md}.
\end{enumerate}

\subsection{Pre-processing for Web Demonstration}
To achieve real-time interactivity within the web-based demonstration, a decoupled inference strategy is employed. This requires migrating the model from a heavy Python environment to a lightweight client-side runtime via the following pre-processing steps:

\subsubsection{ONNX Model Export}
The SAM mask decoder must be exported to the ONNX format to optimize inference performance in the browser. This can be accomplished through two methods:
\begin{itemize}
	\item \textbf{Interactive Export:} Using the\\ \texttt{Source\_code/notebooks/onnx\_model\_example.ipynb} notebook for a guided, step-by-step export process.
	\item \textbf{Script-Based Export:} Executing the automated script from the \texttt{Source\_code} directory:
	\begin{lstlisting}
		python scripts/export_onnx_model.py --checkpoint <path/to/checkpoint> --model-type <model_type> --output <path/to/output>
	\end{lstlisting}
\end{itemize}

\subsubsection{Image Embedding Extraction}
Since the Vision Transformer (ViT) image encoder is computationally expensive, the web demo does not run it locally in the browser. Instead, users must pre-calculate the high-dimensional image embeddings for each target image. This extraction is performed using the utilities provided in \texttt{Source\_code/notebooks/onnx\_model\_example.ipynb}.

\subsubsection{Deployment to Web Frontend}
For the static web application to function, the generated outputs must be manually integrated into the frontend assets. Users must copy the exported ONNX model (\texttt{.onnx}) and the pre-computed image embeddings (\texttt{.npy}) into the \texttt{demo} subfolder. This architecture allows the web app to perform a simple "lookup" of image features, enabling the SAM mask decoder to generate segments in response to user clicks with millisecond latency.

\section{Model Architecture}
The Segment Anything Model (SAM) inference pipeline consists of three primary stages: image encoding, prompt encoding, and mask decoding. This section details the architectural components implemented in the core codebase.

\subsection{Common Architectural Blocks - common.py}
The model utilizes foundational layers defined in \texttt{common.py} to maintain consistency across its sub-networks:
\begin{itemize}
	\item \textbf{MLPBlock:} A standard Feed-Forward Network (FFN) that utilizes two linear layers with a GELU activation function. It projects input embeddings into a higher-dimensional space (\texttt{mlp\_dim}) before restoring the original embedding dimension.
	\item \textbf{LayerNorm2d:} A specialized normalization layer designed for 2D spatial data. Unlike standard \texttt{LayerNorm}, which typically operates on the last dimension, this implementation calculates mean and variance across the channel dimension while preserving spatial height and width. This is essential for maintaining feature map integrity in convolutional "neck" modules.
\end{itemize}

\subsection{Image Encoder - image\_encoder.py}
The \texttt{ImageEncoderViT} serves as the backbone of the system, transforming raw imagery into high-dimensional feature embeddings. The architectural flow and corresponding tensor transformations are summarized below:

\begin{enumerate}
	\item \textbf{Initial Projection:} The input image $x: [B, C_{in}, H_{img}, W_{img}] $ (e.g., $[B, 3, 1024,1024]$) is processed by the \texttt{PatchEmbed} module. This utilizes a convolutional layer with a $16 \times 16$ kernel and stride, divides the image into patches and an convolutional layer is applied, yielding a tensor of shape $[B, H_{p}, W_{p}, D_{embed}]$ (e.g., $[B, 64, 64, 768]$).
%	\item \textbf{Positional Encoding:} If enabled, absolute positional embeddings are added to the patch tokens to provide the transformer blocks with spatial context.
	\item \textbf{Positional Encoding:} To provide the transformer blocks with spatial context, the model integrates absolute positional embeddings. When \texttt{use\_abs\_pos} is enabled, the parameter \texttt{self.pos\_embed} is initialized as a learnable tensor of zeros with the shape $[1, H_{p}, W_{p}, D_{embed}]$ (e.g., $[1, 64, 64, 768]$). During the forward pass, this spatial information is injected into the patch tokens via element-wise addition: $x = x + \text{self.pos\_embed}$. 
	
	This zero-initialization strategy is a standard practice in Vision Transformer architectures, such as ViTDet, to avoid introducing stochastic noise during the initial stages of training. Since the self-attention mechanism is inherently permutation-invariant, it treats patches as an unordered set. By utilizing a learnable parameter starting from zero, the model can adaptively learn optimal spatial encodings tailored specifically for dense segmentation tasks without an initial positional bias. Furthermore, this parameter acts as a placeholder that is typically overwritten by learned values when loading pre-trained checkpoints from large-scale datasets like SA-1B.
	\item \textbf{Transformer Processing:} The data passes through a sequence of \textbf{12 \texttt{Block} modules}, where each block functions as a transformer encoder. Within these blocks:
	\begin{itemize}
		\item Each encoder utilizes \textbf{12 multi-attention heads} to process the embeddings.
		\item The attention mechanism can be configured for either global attention or window-based attention to optimize computational efficiency.
		\item Each block applies \texttt{LayerNorm} before the attention layer and again before a \texttt{MLPBlock}, utilizing residual connections to facilitate gradient flow.
	\end{itemize}
	\item \textbf{Neck and Final Embedding:} To prepare the features for the mask decoder, the tensor is permuted to a channel-first format $[B, D_{embed}, H_{p}, W_{p}]$. It then passes through a "neck" consisting of a $1\times1$ convolution, \texttt{LayerNorm2d}, a $3\times3$ convolution, and a final \texttt{LayerNorm2d} to refine the feature maps.
	\item \textbf{Output:} The final image embedding is produced with shape $[B, C_{out}, H_{p}, W_{p}]$ (e.g., $[B, 256, 64, 64]$), representing the extracted features ready for downstream prompting.
\end{enumerate}